\begin{abstract}
\vspace{-0.3cm}
Test-time scaling for complex reasoning tasks shows that leveraging inference-time compute, by methods such as independently sampling and aggregating multiple solutions, results in significantly better task outcomes. However, a critical bottleneck is \textit{verification}: sampling is only effective if correct solutions can be reliably identified among candidates. While existing approaches typically evaluate candidates independently via scalar scoring, we demonstrate that models are substantially stronger at \textbf{pairwise self-verification}. Leveraging this insight, we introduce \textbf{\method{}}, a framework that unifies generation and verification through efficient pairwise ranking. \method{} comprises two components: \textbf{\method{}-Infer}, an uncertainty-guided algorithm using a tournament-based ranking that dynamically allocates self-verification compute to candidate pairs whose relative correctness is most uncertain; and \textbf{\method{}-PairRL}, an RL framework that \textbf{jointly trains} a single model as both generator and pairwise self-verifier, ensuring the verifier adapts to the generator's evolving distribution. On code generation (LiveCodeBench, CodeContests, SWE-Bench) and math reasoning (AIME, HMMT) benchmarks, \method{}-Infer improves Pass@1 by up to $10\%$ over pointwise verification and outperforms recent test-time scaling methods while being significantly more efficient. Furthermore, \method{}-PairRL achieves $7$--$9\%$ test-time scaling gains over standard RL and pointwise joint training, and improves base Pass@1 by up to 8.7\% over standard RL in a code-generation setting.
\end{abstract}
